\documentclass[11pt]{article}

\usepackage[T1]{fontenc}
\usepackage{newpxtext}
\usepackage{amsmath,amsthm,mathtools}
\usepackage{newpxmath}
\usepackage[
  activate={true,nocompatibility},
  tracking=true,
  expansion=true,
  factor=1100,
  stretch=10,
  shrink=10
]{microtype}
\microtypecontext{spacing=nonfrench}

\usepackage[margin=1.05in]{geometry}
\usepackage{graphicx}
\usepackage[dvipsnames]{xcolor}
\usepackage{booktabs}
\usepackage{array}
\usepackage{enumitem}
\usepackage{caption}
\usepackage{fancyhdr}
\usepackage{titlesec}
\usepackage{needspace}
\usepackage{etoolbox}
\usepackage[most]{tcolorbox}
\usepackage{hyperref}

\hypersetup{
  colorlinks=true,
  linkcolor=linkcolor,
  citecolor=citecolor,
  urlcolor=linkcolor,
  pdfauthor={Lev Petersen},
  pdftitle={Microprice as a Centered Poisson Corrector}
}

\widowpenalty=10000
\clubpenalty=10000
\raggedbottom
\emergencystretch=2em

\definecolor{sectioncolor}{RGB}{18,50,64}
\definecolor{linkcolor}{RGB}{38,80,112}
\definecolor{citecolor}{RGB}{141,93,52}
\definecolor{theorembar}{RGB}{15,118,110}
\definecolor{defnbg}{RGB}{245,248,248}
\definecolor{remarkbar}{RGB}{128,116,143}
\definecolor{titleline}{RGB}{181,136,78}

\titleformat{\section}
  {\normalfont\Large\bfseries\color{sectioncolor}}
  {\thesection}{0.8em}{}
\titleformat{\subsection}
  {\normalfont\large\bfseries\color{sectioncolor}}
  {\thesubsection}{0.8em}{}

\titlespacing*{\section}{0pt}{2.8ex plus 1.0ex minus .6ex}{1.4ex plus .2ex}
\titlespacing*{\subsection}{0pt}{2.0ex plus .8ex minus .4ex}{0.9ex plus .2ex}

\pretocmd{\section}{\needspace{8\baselineskip}}{}{}
\pretocmd{\subsection}{\needspace{6\baselineskip}}{}{}

\pagestyle{fancy}
\fancyhf{}
\fancyhead[L]{\small\itshape Microprice as a Centered Poisson Corrector}
\fancyhead[R]{\small\thepage}
\renewcommand{\headrulewidth}{0.3pt}

\fancypagestyle{plain}{
  \fancyhf{}
  \fancyfoot[C]{\small\thepage}
  \renewcommand{\headrulewidth}{0pt}
}

\theoremstyle{plain}
\newtheorem{proposition}{Proposition}[section]
\newtheorem{corollary}[proposition]{Corollary}
\newtheorem{remark}[proposition]{Remark}

\tcolorboxenvironment{proposition}{
  blanker, borderline west={2.5pt}{0pt}{theorembar},
  left=10pt, before skip=9pt plus 2pt minus 1pt,
  after skip=9pt plus 2pt minus 1pt, breakable
}
\tcolorboxenvironment{corollary}{
  blanker, borderline west={2.5pt}{0pt}{theorembar},
  left=10pt, before skip=9pt plus 2pt minus 1pt,
  after skip=9pt plus 2pt minus 1pt, breakable
}
\tcolorboxenvironment{remark}{
  blanker, borderline west={1.5pt}{0pt}{remarkbar},
  left=10pt, before skip=8pt plus 2pt minus 1pt,
  after skip=8pt plus 2pt minus 1pt, breakable
}

\captionsetup{
  font=small,
  labelfont=bf,
  labelsep=period,
  justification=raggedright,
  singlelinecheck=false,
  skip=5pt,
  hypcap=false
}

\setlength{\abovedisplayskip}{8pt plus 2pt minus 4pt}
\setlength{\belowdisplayskip}{8pt plus 2pt minus 4pt}
\setlength{\abovedisplayshortskip}{6pt plus 2pt minus 3pt}
\setlength{\belowdisplayshortskip}{6pt plus 2pt minus 3pt}

\newcommand{\1}{\mathbf{1}}
\newcommand{\R}{\mathbb{R}}
\newcommand{\E}{\mathbb{E}}
\newcommand{\Prob}{\mathbb{P}}
\newcommand{\norm}[1]{\lVert #1 \rVert}
\newcommand{\PiB}{\Pi_B}

\begin{document}
\thispagestyle{plain}

\begin{center}
{\LARGE\bfseries Microprice as a Centered Poisson Corrector\par}
\vspace{0.35em}
{\large\color{sectioncolor} Paper I: move-time structure, centering, and horizon effects\par}
\vspace{0.7em}
{\large Lev Petersen\par}
\vspace{0.7em}
{\color{titleline}\rule{0.34\textwidth}{0.8pt}}\par
\end{center}

\vspace{0.55em}
\begin{tcolorbox}[
  enhanced,
  colback=defnbg,
  colframe=sectioncolor!18,
  boxrule=0.5pt,
  arc=1.5pt,
  left=8pt,
  right=8pt,
  top=6pt,
  bottom=6pt
]
\textbf{Abstract.} This paper has two goals. First, it identifies the object that classical microprice estimates after the event-time model is reduced to price-move time. For a finite top-of-book Markov model with no-move kernel $Q$, move kernel $T$, event-time one-step drift
\[
r(x):=\E[p_{n+1}-p_n\mid X_n=x],
\]
and fundamental matrix $N=(I-Q)^{-1}$, define
\[
G_1:=Nr,
\qquad
B:=NT.
\]
The $n$-move predictor is then
\[
G^{(n)}=\sum_{j=0}^{n-1}B^jG_1.
\]
Second, it explains why the long-run correction can fail at short horizons even when the move-time model is well specified. If the move-time chain is ergodic with stationary distribution $\pi$, mean move drift $\mu=\pi^\top G_1$, and projector $\PiB:=\1\pi^\top$, then
\[
G^{(n)}=n\mu\,\1+g-B^n g,
\qquad
(I-B)g=G_1-\mu\1,
\qquad
\pi^\top g=0,
\]
and the long-run move-time microprice level is
\[
G^*:=(I-B+\PiB)^{-1}G_1=g+\mu\1.
\]
Thus classical microprice is a centered Poisson corrector on the move-time chain, while the transient term $-B^n g$ governs the horizon effect. We also study the move-time discounted resolvent $(I-\beta B)^{-1}G_1$, where $\beta$ discounts across price moves after the reduction; this is distinct from the event-time discount $\gamma$ used in the companion note \emph{Microprice, Bellman Values, and Semi-Markov Reinforcement Learning}.

Empirically, this decomposition predicts a horizon effect: at short horizons the transient term can still dominate, so $G_1$ or even the raw mid can outperform $G^*$; at longer horizons the centered correction wins when the move-time martingale diagnostic remains small. Public BAC/CVX benchmarks, a five-name WSELOB panel, and a one-day IEX stress test all support that pattern. On WSELOB, the raw mid or $G_1$ is best at 10 seconds, whereas $G^*$ or Disc(.99) is best on all five names by 60 seconds.
\end{tcolorbox}

\section{Introduction}

Stoikov's microprice \cite{Stoikov2018} estimates fair value from top-of-book state by propagating expected future price changes through an absorbing Markov chain. This paper is the structural half of a two-note sequence. Its aim is to identify the long-run object computed by the original construction and to explain the finite-horizon effects that appear when that long-run object is used as a short-horizon forecaster.

Throughout both notes we use the same basic notation. Let
\[
r(x):=\E[p_{n+1}-p_n\mid X_n=x],
\qquad
N:=(I-Q)^{-1},
\qquad
G_1:=Nr,
\qquad
B:=NT.
\]
If the move-time chain $B$ is ergodic with stationary distribution $\pi$, we write
\[
\mu:=\pi^\top G_1,
\qquad
\PiB:=\1\pi^\top,
\qquad
(I-B)g=G_1-\mu\1,
\qquad
\pi^\top g=0,
\]
and
\[
G^*:=(I-B+\PiB)^{-1}G_1=g+\mu\1.
\]
Thus $g$ is the centered move-time relative value, while $G^*$ is the uncentered long-run move-time microprice level. In centered regimes $\mu=0$, so the two coincide.

The division of labor between the two papers is deliberate. The present paper isolates the move-time Poisson structure, the martingale correction, and the empirical horizon crossover. The companion note keeps the same objects $(r,N,G_1,B,g,G^*)$ and adds exact Bellman targets on event time, a semi-Markov interpretation on move time, and online TD-style learning rules. Read together, the first paper explains what microprice is asymptotically, while the second explains how the same construction fits into value-function language and horizon-aware forecasting.

The empirical section evaluates that prediction on three datasets. The public BAC/CVX benchmarks remain useful calibration checks, but they do not show uniform dominance: on BAC the weighted mid is still best under the grouped conditional-curve metric used here. The clearest positive evidence comes from WSELOB, where short- and long-horizon regimes separate cleanly. A one-day IEX venue-local panel then shows the complementary boundary case in which the centering diagnostic fails and discounted extrapolation becomes unreliable.

\section{From event time to move time}

Let $X_n$ be a finite-state Markov approximation of the top of book at event time, and let $p_n$ be the mid-price. Decompose the transition kernel according to whether the mid-price moves:
\[
P=Q+T,
\]
where $Q_{xy}$ is the probability of moving from state $x$ to state $y$ without changing the mid-price and $T_{xy}$ is the probability of moving from $x$ to $y$ with a price change. Define the event-time one-step drift
\[
r(x):=\E[p_{n+1}-p_n\mid X_n=x].
\]

Assume $\rho(Q)<1$, so the no-move subchain is transient and
\[
N:=(I-Q)^{-1}=\sum_{k\ge 0}Q^k
\]
exists. Set
\[
G_1:=Nr,
\qquad
B:=NT.
\]

\begin{proposition}[Embedded move-time chain]\label{prop:embedded}
Let
\[
\tau_0:=0,
\qquad
\tau_{k+1}:=\inf\{n>\tau_k:p_n\neq p_{n-1}\},
\qquad
Y_k:=X_{\tau_k}.
\]
Then $Y_k$ is a Markov chain with transition kernel $B$. Moreover,
\[
G_1(x)=\E_x[p_{\tau_1}-p_{\tau_0}],
\]
so $G_1$ is the expected price change at the next move conditional on the current state.
\end{proposition}

The absorbing-chain recursion therefore reduces to a move-time Markov system: $B$ controls transitions from one price move to the next, and $G_1$ is the corresponding one-move price drift. The same reduction also produces the event-time-discounted objects studied in the companion note: for $\gamma\in(0,1)$,
\[
G_{1,\gamma}:=(I-\gamma Q)^{-1}r,
\qquad
B_\gamma:=\gamma(I-\gamma Q)^{-1}T.
\]
These are the move-time Bellman reward and continuation kernel for the option ``wait until the next mid-price move.'' They discount each no-move excursion in event time before the excursion is collapsed. By contrast, the present paper later studies a move-time resolvent $(I-\beta B)^{-1}G_1$, where the discount is applied only after the reduction to price-move episodes. We therefore reserve $\gamma$ for event-time discounting and $\beta$ for move-time discounting.

If $\pi$ denotes the stationary distribution of $B$, define
\[
\mu:=\pi^\top G_1,
\qquad
\PiB:=\1\pi^\top.
\]

\begin{remark}[Centering]
When $B$ is ergodic, $B^n\to \PiB$. Hence Stoikov's convergence condition $B^nG_1\to 0$ is equivalent to
\[
\pi^\top G_1=0.
\]
Only the invariant constant mode can obstruct convergence.
\end{remark}

\section{The centered Poisson corrector}

Define the $n$-move price adjustment by
\[
G^{(n)}:=\sum_{j=0}^{n-1}B^jG_1.
\]
Equivalently,
\[
\E[p_{\tau_n}-p_{\tau_0}\mid Y_0=x]=G^{(n)}(x).
\]

\begin{proposition}[Finite-horizon decomposition]\label{prop:finite}
Assume that $B$ is ergodic with stationary distribution $\pi$, and set
\[
h:=G_1-\mu\1.
\]
Let $g$ be the unique centered solution of
\[
(I-B)g=h,
\qquad
\pi^\top g=0.
\]
Then, for every $n\ge 1$,
\[
G^{(n)}=n\mu\,\1+g-B^n g.
\]
If
\[
G^*:=(I-B+\PiB)^{-1}G_1,
\]
then
\[
G^*=g+\mu\1.
\]
In particular, if $\mu=0$, then $G^*=g$ and
\[
G^{(n)}\to G^*.
\]
\end{proposition}

The predictor therefore splits into three pieces: a linear drift term $n\mu\1$, a bounded state correction $g$, and a transient remainder $-B^n g$. The centered microprice is the bounded component. The empirical horizon effect comes from the transient term: at short horizons it can still dominate the limiting correction.

\begin{corollary}[Move-time martingale correction]\label{cor:mart}
Under the assumptions of Proposition~\ref{prop:finite},
\[
M_k:=p_{\tau_k}-k\mu+g(Y_k)
\]
is a martingale with respect to the move-time filtration. When $\mu=0$, the corrected price
\[
p_{\tau_k}+G^*(Y_k)
\]
is a move-time martingale.
\end{corollary}

\begin{proposition}[Move-time discounted regularization]\label{prop:disc}
For $\beta\in(0,1)$ define
\[
g_\beta:=(I-\beta B)^{-1}G_1.
\]
Then
\[
g_\beta=\frac{\mu}{1-\beta}\,\1+(I-\beta B)^{-1}(G_1-\mu\1).
\]
Hence the only singular term as $\beta\uparrow 1$ is the constant mode. If $\mu=0$ and, for some $\rho\in(0,1)$,
\[
\norm{B^kG_1}_\infty\le C_\rho\rho^k\norm{G_1}_\infty
\qquad
\text{for all }k\ge 0,
\]
then
\[
\norm{g_\beta-G^*}_\infty
\le
(1-\beta)\frac{C_\rho\rho}{(1-\rho)^2}\norm{G_1}_\infty.
\]
\end{proposition}

The resolvent is always well posed, even when sample centering is imperfect. In centered regimes it is nearly lossless; in drifting or weakly centered regimes its constant mode can dominate. It is therefore best viewed as a regularized proxy and a useful stress test, not as an automatic substitute for the centered correction. It also differs conceptually from the discounted Bellman values in the companion note: $(I-\beta B)^{-1}G_1$ discounts after the reduction to move time, whereas $V_\gamma=(I-\gamma P)^{-1}r$ discounts before the reduction and therefore reweights no-move excursion lengths.

\section{Empirical results}

\subsection{Evaluation protocol}

We evaluate predictors against grouped state-conditioned future-price curves rather than pointwise forecast errors. For a horizon $h$, let $c$ index a spread/imbalance bucket. Write $\bar \Delta_h(c)$ for the average future mid-price change over all rows in bucket $c$, and write $\bar f(c)$ for the average predictor value in the same bucket. The reported score is
\[
\mathrm{RMSE}_h(f):=
\left(
\frac{1}{|C_h|}
\sum_{c\in C_h}\bigl(\bar \Delta_h(c)-\bar f(c)\bigr)^2
\right)^{1/2}.
\]
The horizons are fixed clock/event-time horizons rather than exact move counts, so the empirical question is qualitative: does the long-run correction become more useful as the horizon lengthens? Throughout this paper, ``Disc(.99)'' means the move-time resolvent $(I-0.99\,B)^{-1}G_1$; the companion note reserves $V_{0.99}$ for event-time discounting on $P$.

To connect the data back to Corollary~\ref{cor:mart}, we also report the move-time martingale diagnostic. Writing
\[
\Delta P_k:=p_{\tau_{k+1}}-p_{\tau_k},
\qquad
\widehat \xi_k:=\Delta P_k-\widehat\mu+G^*(Y_{k+1})-G^*(Y_k),
\qquad
\widehat t_M:=\left|\frac{\overline{\widehat \xi}}{s_{\widehat \xi}/\sqrt{n}}\right|.
\]
Small values of $\widehat t_M$ indicate that the estimated correction is close to a move-time martingale.

\subsection{Public BAC/CVX benchmarks}

Table~\ref{tab:public} and Figure~\ref{fig:public} summarize the public BAC/CVX benchmarks. On BAC, $G^*$ reduces grouped RMSE by about $51\%$ at 60 seconds and $76\%$ at 180 seconds relative to the raw mid, but the weighted mid remains best on this metric. On CVX, by contrast, $G^*$ reduces grouped RMSE by roughly $62\%$ and decisively beats the weighted mid. In all three benchmark cases, $\widehat t_M<2$.

\needspace{10\baselineskip}
\par\noindent\begin{minipage}{\textwidth}
\centering
\captionof{table}{Public BAC/CVX grouped-RMSE benchmarks and move-time diagnostics.}
\label{tab:public}
\small
\setlength{\tabcolsep}{4.5pt}
\begin{tabular}{lrrrrrrr}
\toprule
Asset & Horizon (s) & Raw mid & Weighted mid & $G_1$ & $G^*$ & Disc(.99) & $\widehat t_M$ \\
\midrule
BAC & 60  & 0.001812 & 0.000435 & 0.000909 & 0.000892 & 0.000892 & 1.12 \\
BAC & 180 & 0.002287 & 0.000345 & 0.000561 & 0.000547 & 0.000547 & 1.12 \\
CVX & 10  & 0.001828 & 0.003941 & 0.000910 & 0.000697 & 0.000699 & 0.84 \\
\bottomrule
\end{tabular}
\end{minipage}

\vspace{0.65em}
\needspace{17\baselineskip}
\par\noindent\begin{minipage}{\textwidth}
\centering
\includegraphics[width=\textwidth]{figures/public_benchmark_summary.pdf}
\captionof{figure}{Public BAC/CVX benchmarks. Left: grouped RMSE relative to the raw mid. Right: the move-time martingale diagnostic. BAC remains a nuanced benchmark in which the weighted mid is highly competitive, whereas CVX is the cleaner public success case for $G^*$.}
\label{fig:public}
\end{minipage}

\subsection{WSELOB: the horizon crossover}

WSELOB gives the clearest horizon crossover. At 10 seconds, the raw mid is best on PEKAO, PKOBP, and PZU, while $G_1$ is best on KGHM and PKNORLEN. On those same names, $G^*$ is still between $6\%$ and $83\%$ worse than the raw mid, which is exactly what one expects when the transient term $-B^n g$ remains material.

By 60 seconds the ordering reverses. $G^*$ is best on four names and Disc(.99) is best on the fifth. Relative to the raw mid, $G^*$ reduces grouped RMSE by between $20\%$ and $56\%$ on all five names. The weighted mid is not competitive: it is between $4.7\times$ and $11.0\times$ the raw-mid error at 10 seconds, and between $2.1\times$ and $6.6\times$ the $G^*$ error at 60 seconds. The move-time diagnostic remains comfortably in the centered regime, with $\widehat t_M<2$ for every name.

\needspace{11\baselineskip}
\par\noindent\begin{minipage}{\textwidth}
\centering
\captionof{table}{WSELOB summary by asset.}
\label{tab:wselob}
\small
\setlength{\tabcolsep}{4.5pt}
\begin{tabular}{lccrrr}
\toprule
Asset & Best @ 10s & Best @ 60s & Raw$/G^*$ @ 60s & WMid$/G^*$ @ 60s & $\widehat t_M$ \\
\midrule
KGHM      & $G_1$   & $G^*$      & 1.25 & 2.07 & 1.97 \\
PEKAO     & Raw mid & Disc(.99)  & 1.97 & 6.61 & 1.08 \\
PKNORLEN  & $G_1$   & $G^*$      & 1.54 & 4.14 & 0.68 \\
PKOBP     & Raw mid & $G^*$      & 2.03 & 4.94 & 1.85 \\
PZU       & Raw mid & $G^*$      & 2.26 & 6.37 & 0.77 \\
\bottomrule
\end{tabular}
\end{minipage}

\vspace{0.65em}
\needspace{18\baselineskip}
\par\noindent\begin{minipage}{\textwidth}
\centering
\includegraphics[width=\textwidth]{figures/wselob_summary_compact.pdf}
\captionof{figure}{WSELOB summary. Relative RMSE is normalized by the raw mid, so values below one indicate an improvement over the raw mid. The crossover is clear: the limiting correction is not yet best at 10 seconds, but it becomes best by 60 seconds once the horizon is long enough.}
\label{fig:wselob}
\end{minipage}

\subsection{IEX: centering as a diagnostic}

The IEX venue-local panel is deliberately harder: it uses a single venue, a short sample, and a coarse state description, so residual drift is more likely. BAC remains near-centered with $\widehat t_M=0.16$, and $G^*$ is modestly better than the raw or weighted mid at both horizons. SPY and GLD look different. Their diagnostics are $\widehat t_M=3.09$ and $2.24$, respectively, and the discounted surrogate no longer shadows the centered correction. At 10 seconds, Disc(.99) is about $3.16\times$ the $G^*$ error on GLD and $9.50\times$ the $G^*$ error on SPY. This is exactly the failure mode predicted by Proposition~\ref{prop:disc}: once the constant mode is not well centered, discounted extrapolation can be dominated by residual drift rather than by the bounded state correction.

\needspace{18\baselineskip}
\par\noindent\begin{minipage}{\textwidth}
\centering
\includegraphics[width=\textwidth]{figures/iex_summary_compact.pdf}
\captionof{figure}{IEX venue-local stress test. BAC remains near-centered, while SPY and GLD violate the move-time martingale diagnostic. In those names the discounted resolvent no longer behaves like a harmless regularization of $G^*$.}
\label{fig:iex}
\end{minipage}

\section{Conclusion}

Sampling at move times turns the microprice construction into a finite-state Poisson problem. The event-time recursion collapses to a move-time kernel $B$ and a one-move drift $G_1$, while the $n$-move predictor admits the exact decomposition
\[
G^{(n)}=n\mu\,\1+g-B^n g.
\]
This representation clarifies both the martingale correction and the empirical horizon dependence of the estimator.

Across datasets, the centered correction is most useful when two conditions hold at once: the horizon is long enough for the transient term to decay, and the move-time martingale diagnostic remains small. WSELOB is the clearest positive case. BAC/CVX confirm the public benchmark, with the important qualification that BAC still favors the weighted mid under the grouped metric used here. The IEX panel shows the complementary failure mode: when centering weakens, discounted regularization amplifies the wrong mode rather than recovering the bounded state correction.

The companion RL note starts from exactly this dictionary. It keeps $(r,N,G_1,B,g,G^*)$ fixed, replaces the move-time regularizer by exact event-time Bellman values $V^{(H)}$ and $V_\gamma$, and shows that the same price-move episodes support both differential TD for the classical centered correction and semi-Markov TD for discounted forecasting.

Read together, the two papers separate the asymptotic and the objective-matched questions cleanly. The present paper explains the move-time Poisson structure, the martingale diagnostic, and the horizon crossover of $G^*$. The companion note explains how the same reduced model fits into Bellman, average-reward, and semi-Markov RL language without changing the underlying microprice primitives.

\appendix

\section{Proofs}

\begin{proof}[Proof of Proposition~\ref{prop:embedded}]
The first price move occurs after $n$ no-move steps and one move step with kernel $Q^nT$. Summing over $n\ge 0$ gives
\[
\Prob_x(Y_1=y)=\sum_{n\ge 0}(Q^nT)_{xy}
=((I-Q)^{-1}T)_{xy}
=B_{xy}.
\]
Thus $B$ is the move-time transition kernel. For the reward identity,
\[
\E_x[p_{\tau_1}-p_{\tau_0}]
=
\sum_{n\ge 0}(Q^nr)_x
=((I-Q)^{-1}r)_x
=G_1(x).
\]
The Markov property of $(Y_k)$ follows from the strong Markov property at the stopping times $(\tau_k)$.
\end{proof}

\begin{proof}[Proof of Proposition~\ref{prop:finite}]
Write $G_1=\mu\1+h$ with $\pi^\top h=0$. Since $(I-B)g=h$,
\[
\sum_{j=0}^{n-1}B^j h
=
\sum_{j=0}^{n-1}B^j(I-B)g
=
\sum_{j=0}^{n-1}(B^jg-B^{j+1}g)
=
g-B^n g.
\]
Therefore
\[
G^{(n)}
=
\sum_{j=0}^{n-1}B^jG_1
=
n\mu\,\1+g-B^n g.
\]
Next,
\[
(I-B+\PiB)(g+\mu\1)
=
(I-B)g+\PiB g+\mu\PiB\1
=
h+0+\mu\1
=
G_1,
\]
so $G^*=g+\mu\1$. If $\mu=0$, then $G^*=g$ and ergodicity gives
\[
B^ng\to \PiB g=\1\pi^\top g=0,
\]
hence $G^{(n)}\to G^*$. Uniqueness of $g$ under $\pi^\top g=0$ follows because the nullspace of $I-B$ consists of constants.
\end{proof}

\begin{proof}[Proof of Corollary~\ref{cor:mart}]
By Proposition~\ref{prop:finite},
\[
(I-B)g=G_1-\mu\1.
\]
Hence
\[
\E[p_{\tau_{k+1}}-p_{\tau_k}\mid Y_k]=G_1(Y_k),
\]
while
\[
\E[g(Y_{k+1})-g(Y_k)\mid Y_k]=(Bg-g)(Y_k)=-(G_1(Y_k)-\mu).
\]
Adding the two conditional expectations yields
\[
\E[M_{k+1}-M_k\mid Y_k]=0.
\]
\end{proof}

\begin{proof}[Proof of Proposition~\ref{prop:disc}]
Because $\|\beta B\|_\infty\le \beta<1$, the Neumann series converges and
\[
g_\beta=\sum_{k\ge 0}\beta^k B^kG_1.
\]
Using $G_1=\mu\1+h$ and $B^k\1=\1$,
\[
g_\beta
=
\sum_{k\ge 0}\beta^k(\mu\1+B^kh)
=
\frac{\mu}{1-\beta}\,\1+(I-\beta B)^{-1}h.
\]
If $\mu=0$, then $h=G_1$ and
\[
G^*-g_\beta
=
\sum_{k\ge 0}(1-\beta^k)B^kG_1.
\]
The $k=0$ term vanishes. For $k\ge 1$, $1-\beta^k\le k(1-\beta)$, so
\[
\norm{G^*-g_\beta}_\infty
\le
(1-\beta)\sum_{k\ge 1}k\,\norm{B^kG_1}_\infty
\le
(1-\beta)C_\rho\norm{G_1}_\infty\sum_{k\ge 1}k\rho^k.
\]
Since $\sum_{k\ge 1}k\rho^k=\rho/(1-\rho)^2$, the stated bound follows.
\end{proof}

\begin{thebibliography}{99}

\bibitem{Stoikov2018}
S. Stoikov,
\emph{The Micro-Price: A High Frequency Estimator of Future Prices},
Quantitative Finance, 18(12):1959--1966, 2018.

\bibitem{KemenySnell1976}
J. G. Kemeny and J. L. Snell,
\emph{Finite Markov Chains},
Springer, 1976.

\bibitem{Meyer1975}
C. D. Meyer,
The role of the group generalized inverse in the theory of finite Markov chains,
\emph{SIAM Review}, 17(3):443--464, 1975.

\end{thebibliography}

\end{document}
